\Subhead{3.9.1}{Борновское приближение. Метод плавных возмущений}

Пусть $\varepsilon_r=\overline\varepsilon_r+\delta\varepsilon(\bm r)$, фон однороден. Из $(\Delta+k^2)U=-k_0^2\delta\varepsilon U$ и исходящей функции Грина следует интегральное уравнение
\begin{equation}
 U=U_0+k_0^2\int G(\bm r,\bm r')\delta\varepsilon(\bm r')U(\bm r')\dd V'.
\end{equation}
Первое приближение Борна заменяет $U$ под интегралом на невозмущённое $U_0$. При плоском падении в дальней зоне амплитуда рассеяния пропорциональна
\begin{equation}
 f(\bm k_s,\bm k_i)=\frac{k_0^2}{4\pi}\int\delta\varepsilon(\bm r')
 \e^{-\ii(\bm k_s-\bm k_i)\cdot\bm r'}\dd V'.
\end{equation}
Для статистически однородного случайного вещества средняя индикатриса связана со спектром пространственной корреляции $\delta\varepsilon$ в точке переданного волнового вектора $\bm q=\bm k_s-\bm k_i$. Мелкие неоднородности дают широкий угловой спектр, крупные — преимущественно малые углы. Нужна слабость многократного рассеяния и обратного воздействия на падающее поле; одной малости локального $\delta\varepsilon$ на очень длинной трассе может быть недостаточно.

Метод плавных возмущений записывает $U=U_0\exp\psi$, где $\psi$ содержит логарифм амплитуды и фазу. При $U_0\ne0$ подстановка и деление на $U_0\e^\psi$ дают точное равенство
\begin{equation}\label{eq:rytovexact}
 \Delta\psi+2\nabla\ln U_0\cdot\nabla\psi+(\nabla\psi)^2
 =-k_0^2\delta\varepsilon.
\end{equation}
В первом приближении отбрасывают последний, квадратичный по возмущению член. Для $U_0=\e^{\ii kz}$:
\begin{equation}\label{eq:rytovfirst}
 (\Delta+2\ii k\partial_z)\psi_1=-k_0^2\delta\varepsilon.
\end{equation}
В первом порядке $\psi_1=U_1/U_0$, где $U_1$ --- борновская поправка; представление $U=U_0\e^{\psi_1}$ иначе учитывает накопление фазы. Это воспроизводимый вывод из исходного уравнения, не цитата из недоступной монографии. Метод удобен для плавных неоднородностей и малых углов; может описывать значительный накопленный фазовый набег при слабых локальных изменениях. Сильные мерцания, нули поля и каустики ограничивают простую логарифмическую аппроксимацию.
